\section{Project Background}
\label{sec:1-1}

Online shops receive thousands of product reviews every day. Reading them by hand is not feasible, yet the sentiment they carry tells a seller which products disappoint buyers and which are praised.

Turning that text into a signal that can be counted is difficult because a star rating and the words of a review do not always agree. A review titled \emph{``Not as Advertised''} that begins ``Product arrived labeled as Jumbo Salted Peanuts\ldots the peanuts were actually small sized unsalted'' is negative without containing a single strongly negative word. A classifier has to read the sentence, not count words.

This project uses the Amazon Fine Food Reviews corpus~\cite{mcauley2013,kaggle_fine_food}: 568{,}454 reviews posted between October 1999 and October 2012, each with a star rating. The ratings give labels without manual annotation, but they are skewed: after cleaning, about 78\% of the reviews are positive and only 7.6\% are neutral. A classifier that always answers \emph{positive} reaches about 78\% accuracy while recognising nothing about the other classes. Deep learning models that read the text in context are therefore tested on a problem where the obvious shortcut is already strong.

\section{Problem Statement}
\label{sec:1-2}

\begin{itemize}
    \item \textbf{Difficulty in reading sentiment from text.} Sentiment depends on negation (``not bad''), contrast (``good but\ldots'') and context, so counting words is not enough. Models must use word order and context.
    \item \textbf{Imbalanced and approximate labels.} The labels are derived from star ratings, not from human annotation, and the neutral class (3 stars) is rare and ambiguous. A model can look good on accuracy while failing on the minority classes.
\end{itemize}

\section{Research Questions}
\label{sec:1-3}

\begin{enumerate}
    \item How can raw reviews and star ratings be turned into a clean labelled corpus in which the same text never appears in both training and test data?
    \item How skewed are the classes and how long are the reviews, and what do these facts imply for the metric, the loss and the sequence length?
    \item How much does a pre-trained Transformer fine-tuned on this data (BERT) improve over a recurrent network trained from scratch (BiLSTM with attention)?
    \item How far do the two trained models improve on the rule that always answers \emph{positive}, and on which classes?
    \item Where do the models fail, in particular on the neutral class and on reviews whose text contradicts their rating?
\end{enumerate}

\section{Project Objectives}
\label{sec:1-4}

To build and compare deep learning classifiers that assign one of three sentiment labels (negative, neutral, positive) to the text of a food review, using a clean, leakage-safe dataset and an evaluation that is not fooled by the dominant positive class.

\section{Scope and Limitations}
\label{sec:1-5}

\textbf{In scope:} three-class sentiment classification of English review text, data cleaning and splitting, a BiLSTM with attention trained from scratch, a fine-tuned BERT encoder, and the always-positive rule as reference.

\textbf{Out of scope / not claimed:}
\begin{itemize}
    \item Human-annotated sentiment -- labels come from star ratings, so some are wrong by construction (Section~\ref{sec:2-2}).
    \item Predicting the sentiment of \emph{future} reviews -- the split is random, not chronological (Section~\ref{sec:5-8}).
    \item Aspect-level sentiment (which feature of the product is praised or criticised), other languages and other product domains.
    \item A deployed service.
\end{itemize}

\section{Report Structure}
\label{sec:1-6}

The report is organised into chapters that follow the order in which the work was built. The chapters are:

\begin{enumerate}
    \item \textbf{Introduction} -- background, problem statement, research questions, objectives and scope.
    \item \textbf{Problem Statement and Approach} (Chapter~\ref{ch:problem-approach}) -- weak labels, success criteria, metrics and the analytical approach.
    \item \textbf{Dataset and Data Specification} (Chapter~\ref{ch:dataset}) -- the Amazon Fine Food Reviews corpus, its schema, conventions and limitations.
    \item \textbf{Data Understanding} (Chapter~\ref{ch:data-understanding}) -- class distribution, text length, change over time, users, noise and vocabulary, and the design decisions they imply.
    \item \textbf{Data Preparation} (Chapter~\ref{ch:data-preparation}) -- cleaning, model input text, stratified split, fixed test subset and class weights.
    \item \textbf{Modeling} (Chapter~\ref{ch:modeling}) -- the BiLSTM with attention and the fine-tuned BERT.
    \item \textbf{Evaluation} (Chapter~\ref{ch:evaluation}) -- test results, confidence intervals, the neutral class, review length, threshold calibration, errors and limitations.
    \item \textbf{Conclusion} (Chapter~\ref{ch:conclusion}) -- summary, limitations, future work and reproducibility.
\end{enumerate}
